% Sample deck about lectos itself. Build: ./build.sh   Present: lectos slides.pdf slides.tex
% The \note{...} commands are the lecture notes; lectos reads them straight from this file.
\documentclass[aspectratio=169]{beamer}
\usetheme{default}
\usecolortheme{seahorse}
\setbeamertemplate{navigation symbols}{}
\usepackage{amsmath,amssymb}

\title{lectos}
\subtitle{A presenter console for PDF slides}
\author{Sample deck}
\date{}

\begin{document}

\begin{frame}
  \titlepage
  \note{Welcome. This deck is about the tool that is presenting it.
  The slides are a plain PDF; the notes you are reading on the laptop are the
  \texttt{\textbackslash note} commands in the Beamer source.}
\end{frame}

\begin{frame}{What lectos is}
  \begin{itemize}
    \item Projector window for the audience
    \item Presenter console for you
    \item Slides: any tool that makes a PDF
    \item Notes: Markdown or Beamer \texttt{\textbackslash note}
  \end{itemize}
  \note{Two windows, one process, all GTK4.
  \begin{itemize}
    \item The projector shows only the slide.
    \item The console shows the current slide, what the audience sees, and the next slide.
  \end{itemize}
  It does \textbf{not} compile anything: build your PDF the way you always do.}
\end{frame}

\begin{frame}{Notes with mathematics}
  Conservation law:
  \[
    \frac{\partial \mathbf U}{\partial t} + \nabla\cdot\mathbf F(\mathbf U) = \mathbf S(\mathbf U)
  \]
  \note{Equations in the notes are typeset with LaTeX and shown as images, for example the
  governing equation $\partial_t U + \nabla\cdot F(U) = S(U)$ on this slide, or the Roe flux
  \[
    \hat F = \frac{F_L+F_R}{2} - \frac12 \lvert\tilde A\rvert\,(U_R-U_L).
  \]
  If a formula does not compile it stays visible as raw TeX.}
\end{frame}

\begin{frame}{Moving around}
  \begin{itemize}
    \item \texttt{l} / \texttt{h}: next / previous (\texttt{5l})
    \item \texttt{gg}, \texttt{G}, \texttt{12G}, \texttt{:12}: go to slide
    \item \texttt{o}: overview, \texttt{/}: search
    \item \texttt{z}: freeze, \texttt{b}: blackout
    \item \texttt{f} / \texttt{F}: fullscreen
  \end{itemize}
  \note{Vim keys are optional; arrows and Space always work.
  Try \texttt{o} for the thumbnail overview, then click a slide.
  Search \texttt{/} looks through slide text and notes: try ``Roe''.}
\end{frame}

\begin{frame}{Freeze and blackout}
  \begin{itemize}
    \item \textbf{Freeze}: browse privately, the audience keeps the current slide
    \item \textbf{Blackout}: projector goes black
    \item Unfreezing catches the projector up
  \end{itemize}
  \note{Use freeze to look ahead or jump back to check something without the audience
  seeing it. When you unfreeze, the projector jumps to the slide you are viewing.
  Search hits in the notes preview their slide but never move the projector.}
\end{frame}

\begin{frame}{Timer and recovery}
  \begin{itemize}
    \item Elapsed, time left, estimated time at current pace
    \item Page, timer, notes position and layout are saved
    \item Closed the window by accident? Start again and carry on
  \end{itemize}
  \note{The pace estimate is elapsed time divided by slides shown, times slides remaining.
  State lives in one small JSON file under \texttt{\textasciitilde/.local/state/lectos},
  written atomically every couple of seconds.}
\end{frame}

\begin{frame}{Draw on the slide}
  Mouse or pen on the projector window.
  \begin{itemize}
    \item \texttt{c} clears ink on the current slide
    \item Pen pressure changes the stroke width
  \end{itemize}
  \note{Ink is per slide and only kept in memory for now. A laser pointer, persistent ink and a
  chalkboard are on the roadmap.}
\end{frame}

\begin{frame}{Try it}
  \texttt{lectos slides.pdf slides.tex}\par\medskip
  \texttt{lectos slides.pdf notes.md}
  \note{The same deck works with Markdown notes: \texttt{notes.md} uses
  \texttt{<!-- slide: N -->} markers to tie sections to slides.
  Thank you.}
\end{frame}

\end{document}
